% SOURCE_AUTHORITY: sga5_fr_workpass.tex, lines 8452--8467 (LF-normalized unit).
% SOURCE_SHA256: 791F4EFFC5E02832D5D77ED03518C8156D6F07E4C8238B03545DB93D883FBB28
\subsubsection*{1.3. Operações \(\otimes\) e \(\underline{\Hom}\) sobre os \(\mathbb Z_\ell\)-feixes}

\paragraph{1.3.1. Produto tensorial.}
Sejam \(F\) e \(G\) dois objetos da categoria \(\underline{\Hom}(\mathbb N^\circ,\Ab(X))\). Se
\[
F=(F_n)_{n\in\mathbb N}\quad\text{e}\quad G=(G_n)_{n\in\mathbb N},
\]
define-se o produto tensorial \(F\otimes G\) de \(F\) e \(G\) pela fórmula seguinte sobre os objetos e de maneira evidente sobre as flechas:
\[
(F\otimes G)_n=F_n\otimes G_n.
\]
Mais precisamente, define-se claramente um bifuntor aditivo exato à direita, chamado produto tensorial:
\[
\otimes:\Hom(\mathbb N^\circ,\Ab(X))\times \Hom(\mathbb N^\circ,\Ab(X))
\longrightarrow \Hom(\mathbb N^\circ,\Ab(X)).
\]
